\documentclass[11pt]{article}
\usepackage[left=0.82in,right=0.82in,top=0.82in,bottom=1.0in]{geometry}\usepackage{amsmath,amssymb,mathtools,booktabs,microtype,fancyhdr}
\usepackage[dvipsnames]{xcolor}\usepackage[colorlinks=true,linkcolor=NavyBlue,urlcolor=BrickRed]{hyperref}
\definecolor{ink}{HTML}{17212B}\definecolor{accent}{HTML}{315B7D}\definecolor{pale}{HTML}{EDF4F8}\color{ink}
\pagestyle{fancy}\fancyhf{}\lhead{Schur functors of $O(n)$}\rhead{Proof and verification}\cfoot{\thepage}\setlength{\headheight}{14pt}\setlength{\parindent}{0pt}\setlength{\parskip}{5pt}
\newcommand{\Sym}{\operatorname{Sym}}\newcommand{\Tr}{\operatorname{Tr}}\newcommand{\Exp}{\operatorname{Exp}_{\sigma}}
\newcommand{\summarybox}[1]{\begin{center}\fcolorbox{accent}{pale}{\parbox{0.91\linewidth}{#1}}\end{center}}
\title{\textbf{Schur-functor representations of $O(n)$}\\[3pt]\large Plethystic pushforward, polynomial-Gaussian traces, and matrix verification}
\author{Proof and verification note}\date{17 July 2026}
\begin{document}\maketitle
\summarybox{\textbf{Outcome.} The orthogonal formula is correct as a coefficientwise stable plethystic pushforward. The finite trace identity is exact before taking a limit. Five representative $(n,\lambda)$ cases give 40 direct comparisons with maximum error $5.33\times10^{-15}$. Haar $O(9)$ tests recover the invariant quadratic form in $\Sym^2W$, its absence from $\Lambda^2W$, and the odd-degree parity obstruction.}

\section{Stable sigma-MGF and its proof}
Let $W_n=\mathbb C^n$ with the defining action of the compact group $O(n)$, fix $\lambda\vdash d$, and put $V_{\lambda,n}=S_\lambda(W_n)$. For a partition $\tau$,
\begin{equation}\label{eq:coeff}
 [m_\tau]\mathcal S^O_{\lambda,n}
 =\dim\left(\bigotimes_i\Sym^{\tau_i}S_\lambda(W_n)\right)^{O(n)}
 =\mu_{O(n)}\bigl(h_\tau[s_\lambda]\bigr).
\end{equation}
The second equality is simply associativity of plethysm:
$h_\tau[s_\lambda[X]]=(h_\tau[s_\lambda])[X]$.
The defining-representation stable sigma-MGF is $\Exp(h_2)$. Applying its coefficient functional to the fixed symmetric function $h_\tau[s_\lambda]$ proves
\begin{equation}\label{eq:push}
 \boxed{[m_\tau]\mathcal S^O_{\lambda,\infty}
 =\left\langle\Exp(h_2),h_\tau[s_\lambda]\right\rangle.}
\end{equation}
No new limiting invariant theory is needed after the defining case. A conservative sufficient stable range is $n\ge d|\tau|$: the relevant expression has tensor degree $d|\tau|$, where Brauer pairings are independent in that range.

Since $-I\in O(n)$ acts on $S_\lambda(W_n)$ by $(-1)^d$, it acts on the space in \eqref{eq:coeff} by $(-1)^{d|\tau|}$. Therefore
\begin{equation}\label{eq:parity}
 \boxed{[m_\tau]\mathcal S^O_{\lambda,n}=0\quad\text{if }d|\tau|\text{ is odd}.}
\end{equation}
This parity conclusion is exact at finite $n$.

\section{Exact traces and the limiting law}
For every $g\in O(n)$ and $r\ge1$, functoriality gives $S_\lambda(g)^r=S_\lambda(g^r)$, whence
\begin{equation}\label{eq:trace}
 \boxed{\Tr(S_\lambda(g)^r)=s_\lambda[X(g^r)]
 =\sum_{\alpha\vdash d}\frac{\chi^\lambda(\alpha)}{z_\alpha}
 \prod_{a\in\alpha}\Tr(g^{ra}).}
\end{equation}
This formula is exact for every allowed $n$; stability is needed only when replacing traces by Gaussian variables.

The joint moment convergence of finitely many orthogonal trace powers is an
external classical-group trace theorem (Diaconis--Evans, with the
normalization below).  Let independent $A_k\sim N(0,1)$ and set
\begin{equation}
 Z_k^O=\sqrt{k}\,A_k+\mathbf1_{2\mid k}.
\end{equation}
Joint moment convergence of fixed orthogonal trace powers and polynomial substitution in \eqref{eq:trace} give
\begin{equation}
 \Tr(S_\lambda(g)^r)\longrightarrow Q^O_{\lambda,r}
 =\sum_{\alpha\vdash d}\frac{\chi^\lambda(\alpha)}{z_\alpha}\prod_{a\in\alpha}Z_{ra}^O.
\end{equation}
For $\lambda=(2)$, $Q^O_{(2),1}=\tfrac12((Z_1^O)^2+Z_2^O)$ has mean $1$, recording the invariant quadratic form. For $\lambda=(1,1)$, the minus sign makes the mean zero.

\section{Matrix construction and formula cross-check}
The verification draws seeded Haar orthogonal matrices by real Gaussian QR with sign correction. A standard Young symmetrizer $c_\lambda$ is built explicitly on $W_n^{\otimes d}$. If $B$ is an orthonormal basis of $\operatorname{im}c_\lambda$, the direct representation matrix is
\begin{equation}
 D_\lambda(g)=B^*g^{\otimes d}B.
\end{equation}
The formula cross-check never uses this matrix: it computes the defining traces, uses Newton's recurrence $jh_j=\sum_{k=1}^jp_kh_{j-k}$, and evaluates the Jacobi--Trudi determinant $s_\lambda=\det(h_{\lambda_i-i+j})$. Tests compare $\Tr(D_\lambda(g)^r)$ for $r=1,2,3$. They also compare
\begin{equation}
 \prod_i h_{\tau_i}[D_\lambda(g)]
 \quad\text{against}\quad
 \prod_i h_{\tau_i}[s_\lambda][X(g)]
\end{equation}
for $\tau=(),(1),(2),(1,1),(3)$, so the actual sigma-MGF integrands are covered.

\begin{center}\begin{tabular}{@{}llllr@{}}\toprule
$n$ & tested $\lambda$ & checks & purpose & status\\\midrule
3 & $(1),(2),(3)$ & trace and sigma integrands & defining/symmetric powers & PASS\\
4 & $(1,1),(2,1)$ & trace and sigma integrands & alternating/mixed symmetry & PASS\\\midrule
\multicolumn{4}{r}{40 comparisons; maximum absolute error} & $5.33\times10^{-15}$\\\bottomrule
\end{tabular}\end{center}

The seeded Haar diagnostic uses 1,200 matrices in $O(9)$. It gives means $0.9979$ for $\chi_{(2)}$ (target $1$), $-0.0289$ for $\chi_{(1,1)}$ (target $0$), and $0.0116$ for $\chi_{(3)}$ (target $0$). Each lies inside $4.5$ estimated standard errors. Group residuals are at most $3.24\times10^{-15}$.

\section{Reproducibility and interpretation}
Run the module below from the repository root, or open \path{marimo_notebooks/schur_orthogonal.py} with marimo. Exact rows are the proof-driven regression test. Haar rows are seeded numerical diagnostics and are reported with standard errors; they are not presented as proofs of convergence.
\begin{center}\footnotesize\texttt{python -m for\_this\_guy.schur\_functor\_classical\_groups.orthogonal.verification}\end{center}

\section*{References}
S. Howe, \href{https://arxiv.org/abs/2412.19295}{\emph{Random matrix statistics and zeroes of $L$-functions via probability in $\lambda$-rings}}.\\
G. Lehrer and R. Zhang, \href{https://arxiv.org/abs/1207.5889}{\emph{The Brauer Category and Invariant Theory}}.
\\P. Diaconis and S. N. Evans, \emph{Linear functionals of eigenvalues of
random matrices}, Trans. Amer. Math. Soc. 353 (2001), 2615--2633.
\end{document}
